\section{Limitations}

Simulation only. The hardware phase was pre-gated on this evaluation
supporting the mechanism, and the gate closed; every number here is
conditioned on the simulator's service model. The model's constants are
literature-anchored, marked low confidence, and their measured
sensitivity is the paper's principal caveat: the binding constant moves
the representative deficit between 4.4 and 21.5 percent, sign-stable
throughout. Synthetic workloads: no public traces of long-horizon
agentic serving existed at specification freeze; generators, anchors,
and seeds are released, and the adversarial family is constructed, not
observed. Three seeds per cell (Amendment 3): intervals are wider at low
load, and several 0.7x cells have upper bounds crossing zero
individually, though the verdict is invariant to how they are read,
since taking every cell at its most favourable bound yields 33 of 108
reaching the support threshold against the 65 required. Cluster sizes 8
to 32: the 64-device axis was removed for measured infeasibility at the
corrected horizon; the trend across measured sizes is flat, and the
memory-bound regime is probed at one configuration rather than swept.
Value weights are operator-supplied, fixed at registration, and swept in
sensitivity. H2 and H3 are measured on representative configurations,
not crossed with the full grid. The two-cell ablation design cannot
separate any element's effect from its measured seed-noise floor at
three replications; the campaign attributes the mechanism's cost to the
workload's structural duty cycle, and finer per-element attribution
requires swept axes and more seeds (data released). The comparison also
sets a single completion-oriented policy against six request-oriented
ones, so it prices this contract rather than the class of contracts; a
differently designed completion guarantee is untested here. One boundary
is a non-goal rather than a limitation: this layer enforces decisions
about flows, their admission, their bounds and their termination, and
makes no content-level judgement about what a flow says or does. The
service model does not represent chunked prefill, so one mechanism of
the deadline-ordered family is unmodelled.

Four gaps in the specification were surfaced by readers after the paper
was frozen and are named here rather than quietly repaired. The
discipline specifies when a reservation ends, by expiry, exhaustion,
renegotiation or halt, and never specifies what becomes of the
accumulated pinned state on any of those paths, although Section 3 makes
that state load-bearing. The guarantee is scoped to resource-motivated
action and the halt to safety revocation, so an erasure obligation
arriving from outside the serving stack is neither, and the layer has no
correct move when a reservation lifetime exceeds an erasure deadline.
The disposition vocabulary cannot hold a non-terminating state, yet
Section 6.1 measured one: flows admitted, never scheduled, still live at
the horizon, in none of the four categories. And the workload's idle-gap
distribution is uncertain in both directions at once. Deployments
interposing a human confirmation step, as information-flow control
layers for agents do, carry a tail longer than we modelled, which would
deepen the stranding; production measurement of tool-call latency share
\citep{sutradhara2026} suggests our tool-time anchors may instead be long,
which would shorten it. The sensitivity axis contains no arm below its
own anchor. The verdict is robust to the anchor in either direction,
since Section 6.3 measures that shorter gaps produce larger deficits,
but the attribution of the cost to a roughly 5 percent duty cycle is
conditional on it.

Finally, the
registration's one-fix clause means any residual measurement artifact, in
either direction, stands uncorrected in these numbers by design; the
amendment log is the complete record of what was corrected and when.
